\documentclass[withoutpreface,bwprint]{cumcmthesis}

\usepackage{amsmath,amssymb,booktabs,graphicx,tabularx,array,float,url}
% 国赛中文图表题名：显示为“表 1 标题”“图 1 标题”，不使用英文冒号。
\usepackage{caption}
\DeclareCaptionLabelSeparator{cnquad}{\quad}
\captionsetup{labelsep=cnquad}
\captionsetup[table]{name=表}
\captionsetup[figure]{name=图}
\usepackage{gbt7714}
\citestyle{numbers}
\usepackage{hyperref}
\hypersetup{colorlinks=true,allcolors=black,breaklinks=true,hypertexnames=false}

% 优秀论文风格：一级标题居中，显示为“一、问题重述”；二级标题显示为“1.1 问题背景”。
\ctexset{
  section = {format = \centering\heiti\zihao{3}, name = {,、}, number = \chinese{section}, aftername = \quad},
  subsection = {format = \heiti\zihao{-4}, number = \arabic{section}.\arabic{subsection}, aftername = \quad},
  subsubsection = {format = \heiti\zihao{-4}, number = \arabic{section}.\arabic{subsection}.\arabic{subsubsection}, aftername = \quad}
}

\title{基于【核心方法】的【研究对象】建模与优化}
\tihao{A}
\baominghao{000000000000}

\begin{document}

\maketitle

\begin{abstract}
摘要应保持一页以内，采用“问题—方法—结果—结论”的写法。第一段简要说明研究背景、数据来源和总体建模路线，不写空泛口号，不出现真实身份信息。

针对问题一，说明本问的核心任务、所建模型、关键约束和主要结果。正式论文中应给出可核验的定量结果或结果文件名，例如最优方案、误差指标、收益值、覆盖率、命中率等；没有运行结果时只保留占位，不得编造数值。

针对问题二，说明相较问题一新增的不确定因素、约束或目标函数变化，给出求解算法和验证方式。若结果写入附件模板，应在摘要中说明对应结果文件。

针对问题三，说明扩展模型、情景分析、敏感性检验或策略比较结论。最后用一句话概括模型的适用性、稳定性和局限。

\keywords{优化模型\quad 线性规划\quad 蒙特卡洛模拟\quad 敏感性分析}
\end{abstract}

\section{问题重述}
\subsection{问题背景}
简要重构题目背景，说明研究对象、现实约束和建模必要性。不要大段复制题面，应提炼与模型相关的信息。

\subsection{问题要求}
结合附件数据和题目小问，将任务拆解为清晰的输出要求。

\section{问题分析}
\subsection{问题一的分析}
说明问题一的输入、输出、目标函数、约束条件和评价指标，并指出与后续问题的衔接关系。

\subsection{问题二的分析}
说明问题二新增因素如何改变参数、约束或目标函数，以及为什么需要采用相应的优化、预测或仿真方法。

\subsection{问题三的分析}
说明问题三如何在前两问基础上扩展，重点交代数据流、模型流和结果流。

\section{模型假设}
\begin{enumerate}
  \item \textbf{假设 1：} 【假设内容】。解释：【说明该假设对应的公式、约束或数据处理理由，以及可能影响】。
  \item \textbf{假设 2：} 【假设内容】。解释：【说明为何不改变题目核心要求，以及如何在敏感性分析中检验】。
  \item \textbf{假设 3：} 【假设内容】。解释：【说明该假设对求解规模、参数估计或结果解释的作用】。
\end{enumerate}

\section{符号说明}
本文涉及的主要变量符号汇总如下。未在表中声明的局部符号在首次出现处说明。
\begin{table}[H]
  \centering
  \caption{主要符号说明}
  \label{tab:symbols}
  \begin{tabularx}{\textwidth}{cXc}
    \toprule
    符号 & 说明 & 单位 \\
    \midrule
    $i$ & 对象、地块、节点或样本编号 & -- \\
    $t$ & 时间、批次或阶段编号 & -- \\
    $x_{it}$ & 第 $i$ 个对象在阶段 $t$ 的决策变量 & 按题意填写 \\
    $J$ & 目标函数值或综合评价指标 & -- \\
    \bottomrule
  \end{tabularx}
\end{table}

\section{数据预处理}
\subsection{数据来源与字段说明}
说明题目附件、外部数据或自行整理数据的文件名、字段含义、单位、时间/空间粒度和主键。若引用公开数据、软件或文献，应在正文用 \verb|\cite{...}| 标注，并在 \texttt{ref.bib} 中维护对应条目。

\subsection{缺失值与异常值处理}
说明缺失值、重复值、异常值和单位不一致的处理规则。

\subsection{建模参数构造}
说明如何由原始字段计算模型参数、评价指标或约束集合。

\section{模型建立与求解}
\subsection{问题一模型的建立与求解}
先定义决策变量，再给出目标函数和约束。每条约束后应说明其对应的题目条件。
\begin{equation}
  \max Z = \sum_{i}\sum_t r_{it}x_{it} - \sum_i\sum_t c_{it}x_{it}.
\end{equation}
其中 $r_{it}$ 表示收益系数，$c_{it}$ 表示成本系数。

\subsection{问题二模型的建立与求解}
说明在问题一模型上的扩展，包括新增参数、不确定因素、风险度量、分类规则或多目标权衡方式。

\subsection{问题三模型的建立与求解}
说明综合模型、策略生成方法、结果文件输出方式和关键计算步骤。

\section{模型检验}
围绕题目核心结论进行检验，可包括误差分析、残差分析、约束可行性检查、敏感性分析、鲁棒性分析或与基准方案对比。

\section{模型评价与推广}
从适用性、解释性、可复现性、计算复杂度和局限性评价模型。

\bibliographystyle{gbt7714-numerical}
\bibliography{ref}

\appendix
\renewcommand{\thesection}{附录\mbox{\Alph{section}}}
\renewcommand{\thesubsection}{\Alph{section}.\arabic{subsection}}
\ctexset{
  section = {format = \centering\heiti\zihao{3}, name = {,}, number = \thesection, aftername = \quad},
  subsection = {format = \heiti\zihao{-4}, number = \thesubsection, aftername = \quad}
}
\section{支撑材料文件列表}
\begin{table}[H]
\centering
\caption{支撑材料文件列表}
\begin{tabularx}{\textwidth}{lX}
\toprule
文件名 & 功能说明 \\
\midrule
\texttt{code/q1.py} & 问题一模型求解程序，输出对应结果表或结果文件 \\
\texttt{code/q2.py} & 问题二模型求解程序，输出对应结果表或结果文件 \\
\texttt{data\_processed/clean\_data.csv} & 清洗后的建模数据 \\
\texttt{result/} & 论文图表、结果 Excel/CSV 与中间结果 \\
\texttt{ai\_record/AI工具使用详情.pdf} & 若使用 AI 工具，放置支撑材料中的 AI 工具使用详情说明 \\
\bottomrule
\end{tabularx}
\end{table}

\section{源程序代码说明}
附录应包含建模所用的完整、可运行源程序代码，或列出代码文件并在支撑材料中提供。若确实未使用程序，应按官方规范在此明确说明。

\section{AI 工具使用详情说明}
若使用 AI 工具，应按照当年竞赛规则如实说明工具名称、版本或型号、使用环节、关键交互摘要、采纳情况和人工修改情况。若未使用 AI 工具，应按当年规则在参考文献之后或规定位置声明未使用。

\end{document}
